% ECONOMICS_PROBLEM: {"id": "H21-405", "title": "Optimal top-income taxation with avoidance and mobility", "jel_code": "H21", "source_id": "OP-0041"}
% !TeX program = xelatex
\documentclass[11pt,a4paper]{article}
\usepackage[margin=23mm]{geometry}
\usepackage{fontspec}
\setmainfont{Latin Modern Roman}
\usepackage{amsmath,amssymb,mathtools}
\usepackage{xcolor,enumitem,booktabs,longtable,array,xurl,hyperref}
\definecolor{muted}{HTML}{53636C}
\hypersetup{colorlinks=true,urlcolor=blue,linkcolor=blue}
\setlength{\parindent}{0pt}
\setlength{\parskip}{5pt}
\newcommand{\R}{\mathbb R}
\newcommand{\E}{\mathbb E}
\newcommand{\N}{\mathbb N}
\newcommand{\ind}{\mathbf 1}
\newcommand{\Var}{\operatorname{Var}}
\newcommand{\Cov}{\operatorname{Cov}}
\newcommand{\supp}{\operatorname{supp}}
\DeclareMathOperator*{\argmin}{arg\,min}
\DeclareMathOperator*{\argmax}{arg\,max}
\newcommand{\entrymeta}[3]{{\sffamily\small\color{muted}\textbf{\texttt{#1}}\quad|\quad #2\par #3\par}}
\newcommand{\Problemref}[1]{\texttt{#1}}
\newcommand{\JELref}[2]{\texttt{#2} (JEL \texttt{#1})}
\begin{document}
\section*{Optimal top-income taxation with avoidance and mobility}
\textbf{Problem ID:} \texttt{H21-405}\quad \textbf{JEL:} \texttt{H21}\par
% BEGIN ECONOMICS STATEMENT
\label{problem:OP-0041}
\label{atlas:prob:P1}
\noindent\textbf{Original atlas identifier:} P1; entry 41. \textbf{Field:} Public finance and taxation. \textbf{Original tractability label:} Hard.

\noindent\textbf{Classification:} Parameterized theoretical/empirical research agenda; not a certified unsolved theorem.

\subsection*{Definitions and assumptions}
Fix a population of taxpayers with type $\theta$ distributed according to $F$, a top-income threshold $\bar z$, and a baseline policy $p_0$. A policy $p=(\tau,e,a)$ specifies a marginal rate $\tau\in[0,1)$ above $\bar z$, enforcement intensity $e$, and legally permitted income-reclassification rules $a$. Under structural model $m$, agents choose labor, compensation bargaining, reported income, and residence. Let $E_m(p)$ denote the equilibrium set; a prespecified selection $s_m(p)\in E_m(p)$ is required wherever it is non-singleton. Let $EV_\theta(p;m)$ be monetary equivalent variation, measured with a common reference price system, including migration costs and tax payments. Let $R(p;m)$ be revenue from all affected tax bases, $C(p)$ administrative resource cost, $w(\theta)\geq0$ fixed integrable welfare weights, and $\lambda\geq0$ the monetary social value of one unit of net public funds. All expectations are finite.

\subsection*{Formal research question}
\emph{Editorial formulation: the mathematical target below is not a theorem quoted from the references.}

Define the feasible set $\mathcal P$ by legal, reporting, enforcement-capacity, and budget constraints. For a model class $\mathcal M$ consistent with linked taxpayer--firm records and declared identification assumptions, characterize
\[
 V(m)=\sup_{p\in\mathcal P}\left\{\int w(\theta)EV_\theta(p;m)\,dF(\theta)+\lambda[R(p;m)-C(p)]\right\}.
\]
The target is an identified set for welfare differences and $\varepsilon$-optimal top-rate/enforcement combinations for every $m\in\mathcal M$, where $\varepsilon\geq0$ is an explicitly chosen welfare tolerance. Establish which tax or enforcement variation distinguishes real earnings responses, legal shifting, evasion, bargaining transfers, and residence changes. Report when only their combined effect is identified. Identification of a local reported-income elasticity alone does not identify $V(m)$.

\subsection*{Interpretation and scope}
The welfare criterion is a normative primitive, not an empirical consequence of tax records. Revenue must include business and capital taxes displaced by reclassification; otherwise a fiscal transfer between tax bases is mistaken for a social loss. Bargaining-induced income changes may redistribute rents, but their welfare sign depends on who loses them and on welfare weights. Migration may change both the tax base and the chosen population of welfare beneficiaries; retain the fixed initial population unless a different social boundary is expressly justified. Tax-reform comparisons require stated counterfactual trends or exogenous assignment, and instrument-based estimates require exclusion restrictions. The substantive extension is a joint, institution-dependent decomposition with mobility and enforcement, rather than a new claim that standard optimal-top-tax benchmarks are unsolved.

\subsection*{Source audit and status}
[S1] and [S3] contain solved optimal-tax models; [S2] is a policy synthesis. They establish antecedents, not that the expanded problem remains wholly open in October 2026. The atlas provides no explicit theorem or dated claim of unresolved status.

\subsection*{References}
\begin{enumerate}[label={[S\arabic*]},leftmargin=*,itemsep=0.6em]
\item Emmanuel Saez (2001). Using Elasticities to Derive Optimal Income Tax Rates. The Review of Economic Studies 68(1), 205--229. DOI: 10.1111/1467-937X.00166.
\url{https://academic.oup.com/restud/article-abstract/68/1/205/1568609}
\newline\emph{Locator and support limits:} Abstract and article, pp. 205--229: optimal-tax benchmark under specified elasticities; does not identify mobility responses.
\item Peter Diamond and Emmanuel Saez (2011). The Case for a Progressive Tax: From Basic Research to Policy Recommendations. Journal of Economic Perspectives 25(4), 165--190. DOI: 10.1257/jep.25.4.165.
\url{https://www.aeaweb.org/articles?id=10.1257/jep.25.4.165}
\newline\emph{Locator and support limits:} Abstract and policy discussion: synthesis of progressive-tax arguments, not a universal implementable rate.
\item Thomas Piketty, Emmanuel Saez, and Stefanie Stantcheva (2014). Optimal Taxation of Top Labor Incomes: A Tale of Three Elasticities. American Economic Journal: Economic Policy 6(1), 230--271. DOI: 10.1257/pol.6.1.230.
\url{https://pubs.aeaweb.org/doi/10.1257/pol.6.1.230}
\newline\emph{Locator and support limits:} Abstract and model: labor supply, avoidance, and bargaining; international mobility is an additional margin.
\end{enumerate}

\subsection*{Common conventions: Source mathematical conventions}

Unless an entry states otherwise, agent and firm sets are finite. State and action spaces are measurable, strategies and outcome maps are measurable, probability laws are specified, and expectations exist for the targets under discussion. Infinite-horizon discount factors lie in $(0,1)$ and discounted objectives are finite. Norms, tolerances and complexity input sizes must be specified for computational comparisons. Constants or primitives that appear locally are inputs, not empirical facts asserted by this catalogue.

\textbf{Identification.} A statistical model $m\in\mathcal M$ implies an observable law $P_m$ and target $\theta(m)$. At law $P$, its identified set is
\[
\Theta(P;\mathcal M)=\{\theta(m):m\in\mathcal M,\ P_m=P\}.
\]
Point identification holds when this set is a singleton, or equivalently when $P_{m_1}=P_{m_2}$ implies $\theta(m_1)=\theta(m_2)$ for all compatible models. Sharp bounds mean bounds attained, or approached when the boundary is not attained, by compatible models. Writing this set is a definition, not a solution: the substantive task is to establish informative restrictions and derive or estimate the set. An unrestricted-confounding model may give only trivial bounds.

\textbf{Inference.} Identification, estimability and valid uncertainty quantification are separate questions. Fix $\alpha\in(0,1)$. A confidence procedure $C_n$ over a declared law class $\mathcal P$ has asymptotic uniform coverage at level $1-\alpha$ when
\[
\liminf_{n\to\infty}\inf_{P\in\mathcal P}\Pr_P\{\theta(P)\in C_n\}\geq1-\alpha.
\]
For set-identified targets the coverage event must be defined separately, for example coverage of the full identified set. If an entry uses identification-indexed models instead of uniquely indexed laws, the same coverage convention is applied over those models. Pointwise limits do not establish uniform validity, and asymptotic validity does not guarantee finite-sample precision. Sample sizes, dependence structures and weak-identification sequences are part of each application.

\textbf{Counterfactuals.} $Y_i(a)$ is a potential outcome under a specified intervention, including its timing, implementation and versions. It is not $Y_i$ conditional on voluntarily choosing $a$. A full intervention vector permits interference; shorthand scalar treatment notation does not silently impose no interference. Assignment, selection, attrition, anticipation and measurement must be specified. A claim about an instrument's local effect is not automatically a population policy effect.

\textbf{Economic equilibrium.} A counterfactual model includes best responses, constraints, feasibility, market clearing, state transitions and expectations consistent with the stated information. When several equilibria exist, either a declared selector is an input or the target is set-valued. Comparisons across policy regimes must use the stated selection convention. Reduced-form relationships are not presumed invariant after an equilibrium-changing intervention.

\textbf{Optimization and welfare.} Optimization is over the stated feasible set. If attainment is not established, $\sup$ or $\inf$ and $\varepsilon$-optimal choices replace an asserted maximizer or minimizer. For example, with value $V=\sup_{a\in\mathcal A}W(a)<\infty$, an $\varepsilon$-optimal set is $\{a\in\mathcal A:W(a)\geq V-\varepsilon\}$ for $\varepsilon>0$. Benefits, costs, risk and health or environmental outcomes can be aggregated only using declared units and conversion weights. Interpersonal utility weights, the treatment of fiscal transfers and foreign profits, discounting, and the behavioral welfare criterion are explicit inputs. A comparative-static derivative additionally requires existence and appropriate differentiability of the selected counterfactual.

\textbf{Sources.} Local labels [S1], [S2], and so forth restart in every question. Locators identify the relevant abstract, model, section or result at the level actually checked; a bibliographic or abstract-level verification is not represented as inspection of an entire proof. Working papers are distinguished from later publications and changing versions. Source summaries are paraphrased. All URL checks and editorial formulations refer to the audit date stated above.
% END ECONOMICS STATEMENT
\end{document}
